% Options for packages loaded elsewhere
\PassOptionsToPackage{unicode}{hyperref}
\PassOptionsToPackage{hyphens}{url}
\documentclass[
]{book}
\usepackage{xcolor}
\usepackage{amsmath,amssymb}
\setcounter{secnumdepth}{5}
\usepackage{iftex}
\ifPDFTeX
  \usepackage[T1]{fontenc}
  \usepackage[utf8]{inputenc}
  \usepackage{textcomp} % provide euro and other symbols
\else % if luatex or xetex
  \usepackage{unicode-math} % this also loads fontspec
  \defaultfontfeatures{Scale=MatchLowercase}
  \defaultfontfeatures[\rmfamily]{Ligatures=TeX,Scale=1}
\fi
\usepackage{lmodern}
\ifPDFTeX\else
  % xetex/luatex font selection
\fi
% Use upquote if available, for straight quotes in verbatim environments
\IfFileExists{upquote.sty}{\usepackage{upquote}}{}
\IfFileExists{microtype.sty}{% use microtype if available
  \usepackage[]{microtype}
  \UseMicrotypeSet[protrusion]{basicmath} % disable protrusion for tt fonts
}{}
\makeatletter
\@ifundefined{KOMAClassName}{% if non-KOMA class
  \IfFileExists{parskip.sty}{%
    \usepackage{parskip}
  }{% else
    \setlength{\parindent}{0pt}
    \setlength{\parskip}{6pt plus 2pt minus 1pt}}
}{% if KOMA class
  \KOMAoptions{parskip=half}}
\makeatother
\usepackage{longtable,booktabs,array}
\newcounter{none} % for unnumbered tables
\usepackage{calc} % for calculating minipage widths
% Correct order of tables after \paragraph or \subparagraph
\usepackage{etoolbox}
\makeatletter
\patchcmd\longtable{\par}{\if@noskipsec\mbox{}\fi\par}{}{}
\makeatother
% Allow footnotes in longtable head/foot
\IfFileExists{footnotehyper.sty}{\usepackage{footnotehyper}}{\usepackage{footnote}}
\makesavenoteenv{longtable}
\usepackage{graphicx}
\makeatletter
\newsavebox\pandoc@box
\newcommand*\pandocbounded[1]{% scales image to fit in text height/width
  \sbox\pandoc@box{#1}%
  \Gscale@div\@tempa{\textheight}{\dimexpr\ht\pandoc@box+\dp\pandoc@box\relax}%
  \Gscale@div\@tempb{\linewidth}{\wd\pandoc@box}%
  \ifdim\@tempb\p@<\@tempa\p@\let\@tempa\@tempb\fi% select the smaller of both
  \ifdim\@tempa\p@<\p@\scalebox{\@tempa}{\usebox\pandoc@box}%
  \else\usebox{\pandoc@box}%
  \fi%
}
% Set default figure placement to htbp
\def\fps@figure{htbp}
\makeatother
\setlength{\emergencystretch}{3em} % prevent overfull lines
\providecommand{\tightlist}{%
  \setlength{\itemsep}{0pt}\setlength{\parskip}{0pt}}
\usepackage[]{biblatex}
\addbibresource{book.bib}
\usepackage{ctex}

%\usepackage{xltxtra} % XeLaTeX的一些额外符号
% 设置中文字体
%\setCJKmainfont[BoldFont={黑体},ItalicFont={楷体}]{新宋体}

% 设置边距
\usepackage{geometry}
\geometry{%
  left=2.0cm, right=2.0cm, top=3.5cm, bottom=2.5cm} 

\usepackage{amsthm,mathrsfs}
\usepackage{booktabs}
\usepackage{longtable}
\makeatletter
\def\thm@space@setup{%
  \thm@preskip=8pt plus 2pt minus 4pt
  \thm@postskip=\thm@preskip
}
\makeatother
\usepackage{bookmark}
\IfFileExists{xurl.sty}{\usepackage{xurl}}{} % add URL line breaks if available
\urlstyle{same}
% fallback for those not using the hyperref driver hyperxmp:
\makeatletter
\@ifundefined{xmpquote}{\newcommand{\xmpquote}[1]{#1}}{}
\makeatother
\hypersetup{
  pdftitle={《黄帝内针讲义》学习笔记},
  pdfauthor={谢家树博士（香港）},
  hidelinks,
  pdfcreator={LaTeX via pandoc}}

\title{《黄帝内针讲义》学习笔记}
\author{谢家树博士（香港）}
\date{2026-07-04}

\begin{document}
\maketitle

{
\setcounter{tocdepth}{1}
\tableofcontents
}
\chapter*{引言}\label{ux5f15ux8a00}
\addcontentsline{toc}{chapter}{引言}

《黄帝内针讲义》的主編是刘力红，广西首位中医博士，现任广西中医药大学基础医学院教师，经典中医临床研究所首席教授。

「黄帝内针」针法源自《黄帝内经》。特色是``六经辫证、易学高效、
技术安全''，屬于经络讲证针灸学派。《黄帝内针讲义》是医药院校的中医教学教材。

全书分为总论和各论。总论主要介绍了黄帝内针的理论体系和临床法则。各论按照疾病所在部位进行分类，详细论述了诊疗方法、医案示例，并附有练习题。

以下是在下自学的笔记。笔记内容编排是按该书的目录排版。

\chapter{前言}\label{ux524dux8a00}

\section{黄帝内针特色、学派}\label{ux9ec4ux5e1dux5185ux9488ux7279ux8272ux5b66ux6d3e}

黄帝内针是源自《黄帝内经》的针法，以``六经辨证、易学高效、技术安全''为特色，属于经络辨证针灸学派。

\section{针灸学术流派}\label{ux9488ux7078ux5b66ux672fux6d41ux6d3e}

当代常用的针灸学术流派大略分为：

\begin{itemize}
\item
  经络脏腑派：以经典的经络辨证、脏腑辨证、八纲辨证、气血津液辨证等为辨证特色的针灸学派。
\item
  阴阳应象派：以传统的应象理论、河洛易理、五运六气等为辨证特色的针灸学派。
\item
  专病专穴派：以``特效穴位、穴证相对''为特色的针灸学派。
\item
  生理解剖派：虽基于生理解剖而非中医理论，但临床疗效同样卓著的针灸学派。
\end{itemize}

\chapter{黄帝内针的淵源与传承}\label{ux9ec4ux5e1dux5185ux9488ux7684ux6df5ux6e90ux4e0eux4f20ux627f}

\section{黄帝内针的淵源}\label{ux9ec4ux5e1dux5185ux9488ux7684ux6df5ux6e90}

黄帝内针代代皆为单传秘授。

\begin{itemize}
\tightlist
\item
  杨其海先生自父杀杨运濟手中传得针灸一脉。
\item
  2016 年年底，由黄帝内针传人杨真海先生传讲、黄帝内针传承弟子刘力红先生整理的 《黄帝内针一一和平的使者》一书正式公开出版。
\end{itemize}

\section{黄帝内针的传承}\label{ux9ec4ux5e1dux5185ux9488ux7684ux4f20ux627f}

传承的路径大致有三条：

\begin{itemize}
\tightlist
\item
  文字传承、
\item
  口耳传承和
\item
  直接传承。
\end{itemize}

传承要着重把握三个要点：

\begin{itemize}
\tightlist
\item
  信（确信皇帝内针的理法方针）、
\item
  诚（按照孝悌忠信礼义廉耻的要求做好人）和
\item
  明（明针法背后的至深之理）。
\end{itemize}

\section{理论依据}\label{ux7406ux8bbaux4f9dux636e}

从阴引阳，从阳引阴，以右治左，以左治右。

阳病治阴，阴病治阳，定其血气，各守其乡。

\chapter{黄帝内针理法概述}\label{ux9ec4ux5e1dux5185ux9488ux7406ux6cd5ux6982ux8ff0}

\section{三才、三焦}\label{ux4e09ux624dux4e09ux7126}

\subsection{三才}\label{ux4e09ux624d}

天有阴阳，人有阴阳，地也有阴阳，生命是三凑六合而成。

无论从法理还是针道的应用，``三''都显得非常重要。

\subsection{三焦}\label{ux4e09ux7126}

在传统中医上的理解，三焦是指六腑之一，属于手少阳经。三焦之腑与其他五腑有其特殊性，在解剖学上完全找不到相应的部位。

黄帝内针的``三焦''概念，是从另一个角度来看待的三焦，是与``三才''相关联的三焦定位。

大致而言，上焦是心窝鸠尾以上的区域，中焦是鸠尾至肚脐神阙的区域， 下焦是神阙以下的区域。

上焦应天，中焦应人，下焦应地。

\section{阴阳}\label{ux9634ux9633}

\subsection{阴阳的含义}\label{ux9634ux9633ux7684ux542bux4e49}

一切的变化，都是阴阳所生。

定则：中医治病有许多方法，但千法万法都不离阴阳。

``本'' 与``末''是相对的概念，在先者为本，在后者为末，本末的原则实质由 先后来确定。

所 谓 `` 治 病必求于本''，一般的理解就是治病必领求到阴阳的层面。

脏腑本身可分阴阳，脏为阳，腑为阳，是从形质中取判阴阳。

从阴阳来论有无，以阴阳来论形气，那又可说气为阳，形为阴，有为阴，无为阳。

中医这个体系的重要特征：

\begin{itemize}
\tightlist
\item
  从脏腑和阴阳来论，更强调阴阳；
\item
  从形与气来论，更强调气；
\item
  从可见（有）和不可见（无）来论，更强调不可见。
\end{itemize}

中医人，眼中什至可以没有脏腑，但一刻都不能设有阴阳。

\subsection{三阴三阳}\label{ux4e09ux9634ux4e09ux9633}

``六经'' 即三阴三阳。

黄帝内针是六经辩证。

三阴即：

\begin{itemize}
\tightlist
\item
  太阴、
\item
  少阴、
\item
  厥阴；
\end{itemize}

三阳即：

\begin{itemize}
\tightlist
\item
  大阳、
\item
  阳明、
\item
  少阳。
\end{itemize}

\subsection{人体层面}\label{ux4ebaux4f53ux5c42ux9762}

脏腑是生命形态的内核机关。

从人的层面说，三阴三阳涵括了人体的经络系统，即三阴经、三阳经。由于经分手足，既有手三阴经、手三阳经，还有足三阴经、足三阳经。

手三阴经：

\begin{itemize}
\tightlist
\item
  手太阴肺经、
\item
  手少阴心经、
\item
  手厥阴心包经。
\end{itemize}

手三阳经：

\begin{itemize}
\tightlist
\item
  手阳明大肠经、
\item
  手太阳小肠经、
\item
  手少阳三焦经。
\end{itemize}

足三阴经：

\begin{itemize}
\tightlist
\item
  足太阴脾经、
\item
  足少阴肾经、
\item
  足厥阴肝经。
\end{itemize}

足三阳经：

\begin{itemize}
\tightlist
\item
  足阳明胃经、
\item
  足太阳膀胱经、
\item
  足少阳胆经。
\end{itemize}

三阴三阳牵涉两个方面：

\begin{itemize}
\tightlist
\item
  脏腑层面、
\item
  经络层面。
\end{itemize}

经络的作用至少有两重：

\begin{itemize}
\tightlist
\item
  联系个体生命形态的内内外外，
\item
  作为个体生命形态与天地之间的重要交通。
\end{itemize}

\subsection{天地层面}\label{ux5929ux5730ux5c42ux9762}

三阴三阳说的是六气。六气通俗一点的表达就是：

\begin{itemize}
\tightlist
\item
  风、
\item
  寒、
\item
  暑、
\item
  湿、
\item
  燥、
\item
  火。
\end{itemize}

学术一点的表达是：

\begin{itemize}
\tightlist
\item
  风木、
\item
  應水、
\item
  相火、
\item
  君火、
\item
  燥金、
\item
  湿士。
\end{itemize}

三阴：厥阳风木、少阴君火、太阴湿土。

三阳：少阳相火、阳明燥金、太阳寒水。

\section{中和}\label{ux4e2dux548c}

调摄阴阳：``中和''。

阴阳的作用是相对的，之所以能够中 (zhong) 节，之所以能够处和，是因为有``中'' 的作用。中的作用天生就是致和。

由矛盾走向协和，这在中医可称之为``平人''。不病，亦即健康的状态。

取穴下针的唯一目的，就是实现阴阴自和。阴阳自和实际上就是以平为期，有中则能平，有中则有和。

\chapter{黄帝内针的基本特点}\label{ux9ec4ux5e1dux5185ux9488ux7684ux57faux672cux7279ux70b9}

\section{执两用中}\label{ux6267ux4e24ux7528ux4e2d}

执两用中是黄帝内针的全部奥妙。

两就是对，亦是双，故日成双成对。这个两端，实质上亦就是左在、上下，亦就是阴阳。具体而言，比如病表现在左，可视为一端，那么治必须在另一端的右，这才构成了两。因此，以右治左、以左治右才符合执其两端的原则。只有符合了这个原则，才能实现用中的目的。

\section{随证治之}\label{ux968fux8bc1ux6cbbux4e4b}

证是中医的服目，是中医人的下手处。``证'' 与``症'' 的含义不完全相同，``证'' 可以包含〝症'' 的内酒，而``症'' 则未必能包含``证''。

证是患者对身体的综合表达，这个表达既包括了症，也就是疾病的表现，也包括了病因，同时还恩合香对问题所给品的自治方茶。因此，证实际上涵括了病证、病因、病治，是三合一。

证（症）可以表达成是机体能够 受到的异常，而机体常见的证（症）
不外酸、麻、胀、痛、痒、熱、寒等，当然还有二便、饮食、呼吸、睡眠等异常。

证有定位的证，比如身体局部疼痛;也有定性的证，比如恶寨、发熱、失眠等。

\section{同气相求}\label{ux540cux6c14ux76f8ux6c42}

又叫〝求同气''，就是求病证的同气。

同气相求，要在有求必应，同气相求是为了有应。求同气，就是求病证的同气，病证在哪里，治就是取同气，病证在哪一部位，治所取的穴就在哪一部位；病证在哪一经，治所取的穴就在哪一经。

\section{导引}\label{ux5bfcux5f15}

导从心入，所以必须透过感来实现。导引就是透过感来实现身心的贯通，身心能够贯通，自然就形与神俱了。

\section{``病'' 与 ``工'' 的导引}\label{ux75c5-ux4e0e-ux5de5-ux7684ux5bfcux5f15}

导引则是要实现医患的和合。

内针不在乎针感上的得气，针感上得气与否丝老不会影响疗效。

针入以后，医者与患者共同关注的焦点是病患之处，医者以言导引之，患者以
意关注之，病处的变化便当即发生。

内针施治的整个过程不行针，一般留针三刻（45 分钟）后除针。

\section{神与形俱}\label{ux795eux4e0eux5f62ux4ff1}

\subsection{君火与相火}\label{ux541bux706bux4e0eux76f8ux706b}

\subsection{守神与守形}\label{ux5b88ux795eux4e0eux5b88ux5f62}

\subsection{心法与针法}\label{ux5fc3ux6cd5ux4e0eux9488ux6cd5}

\printbibliography

\end{document}
